% Task 7.2 — Introduction (681 words, completed 2026-04-20)
% Source: manuscript/introduction.md

Roughly 850 million people worldwide are estimated to have chronic kidney
disease, and the global prevalence grew by 33\% between 1990 and 2017
\citep{Francis2024}. Among people with diabetes, as many as 1 in 3 are
affected; among those with hypertension in high-income settings, the proportion
reaches approximately 1 in 5 \citep{Francis2024}. By 2040, CKD is projected to
rank fifth among leading causes of years of life lost globally
\citep{Francis2024}. Those numbers create genuine pressure on health systems to
identify high-risk patients early, before irreversible loss of kidney function
closes off the best treatment options.

Machine learning has been proposed as a practical answer to this challenge.
Models trained on electronic health records, biomarker panels, and demographic
variables have reported AUROC values above 0.95 in national cohort studies
\citep{Krishnamurthy2021, Bai2022, Li2025, Sabanayagam2025}. Established risk
equations like the Kidney Failure Risk Equation have been validated across
populations in North America, the United Kingdom, and Latin America, confirming
algorithmic CKD risk prediction is technically achievable
\citep{Tangri2016, Major2019, BravoZuniga2025}. The field has not struggled to
build models. The struggle is with what happens after the model is built.

Discrimination metrics like AUROC measure whether a model ranks patients
correctly relative to each other. They say nothing about whether the assigned
probability scores are trustworthy in absolute terms. A model posting AUROC 0.97
assigns a 65\% risk score to patients whose true event rate sits near 20\%, and
clinicians making treatment decisions from those numbers are working from
miscalibrated information. Van Calster and colleagues identified calibration as
the Achilles heel of predictive analytics, noting poor calibration regularly
persists even when discrimination appears strong \citep{VanCalster2019}. A
systematic review of CKD risk models by Echouffo-Tcheugui and Kengne found
calibration is assessed less commonly than discrimination across the published
literature; of all the models reviewed, only eight for CKD occurrence and five
for CKD progression had been externally validated for calibration
\citep{EchouffoTcheugui2012}. The models exist. The evidence for trusting their
probability outputs largely does not.

The problem runs deeper than calibration alone. Campagner and colleagues
reviewed machine learning studies in healthcare and found fewer than 4\% address
uncertainty quantification explicitly \citep{Campagner2025}. A model outputting
a probability without any indication of how much to trust the number puts
clinicians in a difficult position. Banerji and colleagues put this plainly:
clinical AI tools must communicate predictive uncertainty at the level of the
individual patient, not purely in aggregate performance statistics
\citep{Banerji2023}. A predicted CKD risk of 78\% warrants a different clinical
response when the model's uncertainty is narrow than when high uncertainty
renders the output practically unreliable.

A 2023 systematic review commissioned to support CDC prevention guidelines
reached a pointed conclusion: CKD risk prediction models need to be better
calibrated and externally validated before incorporation into clinical guidelines
\citep{GonzalezRocha2023}. No published study in the CKD literature has
operationalized all three of these demands together. No existing work has jointly
evaluated calibration across multiple post-hoc correction methods, quantified
uncertainty through a coverage-guaranteed conformal framework, and assessed
deployment readiness through a structured multi-criterion checklist, all on the
same model suite and across an independent external cohort.

This study addresses the gap. Using the UCI CKD dataset for model development
and MIMIC-IV as an external validation cohort, we trained five classifiers
spanning the range commonly used in clinical prediction: logistic regression,
random forest, gradient boosting via XGBoost, a support vector machine with
Platt-scaled probabilities, and Gaussian naive Bayes. Each model was evaluated
across three dimensions: calibration before and after post-hoc recalibration
using Platt scaling and isotonic regression; predictive uncertainty through split
conformal prediction with a formal 90\% marginal coverage guarantee; and a
structured eight-criterion deployment readiness framework grounded in current
reporting standards including TRIPOD+AI \citep{Collins2024TRIPOD}.

The study has three objectives:
\begin{enumerate}
    \item Quantify pre- and post-calibration error for five CKD classifiers on
          both the internal UCI test set and the external MIMIC-IV cohort.
    \item Apply split conformal prediction to generate prediction sets with a
          90\% coverage guarantee and determine whether the guarantee holds on
          external data.
    \item Score each model against an eight-criterion deployment readiness
          checklist and identify which, if any, meet the threshold for
          responsible clinical use.
\end{enumerate}
